% ============================================================================
\section{可插拔接口：\texttt{syscall\_ops\_t}}
\label{sec:syscall-interface}
% ============================================================================

与 \texttt{pmm\_ops\_t}（第~\ref{ch:pmm}~章）、\texttt{heap\_ops\_t}（第~\ref{ch:kmalloc}~章）、
\texttt{vma\_ops\_t}（第~\ref{ch:vma}~章）完全相同的"dispatch + 后端"模式。
不同的系统调用入口机制被抽象为后端，共享同一套 dispatch 层和 syscall table。

\codefile{src/include/kernel/syscall.h}
\begin{lstlisting}[style=cstyle]
typedef struct syscall_ops {
    const char* name;       /* "int0x80", "sysenter" */

    /*
     * init -- 初始化后端
     *
     * int 0x80 后端：在 IDT[0x80] 安装 DPL=3 中断门
     * sysenter 后端：配置 IA32_SYSENTER_CS/EIP/ESP MSR
     */
    void (*init)(void);
} syscall_ops_t;
\end{lstlisting}

\begin{figure}[H]
  \centering
  \begin{tikzpicture}[
    box/.style={draw, rounded corners=3pt, minimum width=3.0cm, minimum height=0.8cm,
                font=\small, align=center},
    api/.style={box, fill=green!8},
    dispatch/.style={box, fill=yellow!12},
    backend/.style={box, fill=blue!8},
    table/.style={box, fill=orange!10},
    arrow/.style={-{Stealth[length=4pt]}, thick},
    darrow/.style={-{Stealth[length=4pt]}, thick, dashed},
    note/.style={font=\scriptsize\itshape, text=gray},
  ]

  % 调用者
  \node[api] (caller) at (0, 0) {调用者\\{\footnotesize kmain.c / 后端 stub}};

  % 公开 API
  \node[api] (api) at (0, -1.5) {公开 API\\{\footnotesize \texttt{syscall\_init()}}\\{\footnotesize \texttt{syscall\_dispatch()}}};

  % dispatch 层
  \node[dispatch] (disp) at (0, -3.5) {\texttt{syscall.c}\\{\footnotesize \texttt{g\_syscall\_ops} 指针}};

  % 后端
  \node[backend] (int80) at (-3.0, -5.5) {\texttt{int 0x80}\\{\footnotesize IDT gate DPL=3}};
  \node[backend] (sysenter) at (3.0, -5.5) {\texttt{sysenter}\\{\footnotesize MSR 配置}};

  % syscall table
  \node[table] (table) at (0, -7.5) {\texttt{syscall\_table[256]}\\{\footnotesize nr → handler}};

  % handler
  \node[api] (handlers) at (0, -9.2) {handler 实现\\{\footnotesize \texttt{sys\_write} / \texttt{sys\_exit} / \texttt{sys\_brk}}};

  % 箭头
  \draw[arrow] (caller) -- (api);
  \draw[arrow] (api) -- (disp) node[right, pos=0.5, note] {\texttt{g\_syscall\_ops->init()}};
  \draw[arrow] (disp) -- (int80) node[left, pos=0.5, note] {D-3};
  \draw[darrow] (disp) -- (sysenter) node[right, pos=0.5, note] {D-4（进阶）};

  % 后端调用 dispatch
  \draw[arrow, red!60!black] (int80.south) |- (table.west)
    node[below left, pos=0.25, note] {stub 调用};
  \draw[darrow, red!60!black] (sysenter.south) |- (table.east);

  \draw[arrow] (table) -- (handlers) node[right, pos=0.5, note] {\texttt{table[nr](regs)}};

  % kconfig 标注
  \node[font=\scriptsize, draw, rounded corners, fill=purple!5, text=purple!60!black,
        align=center, text width=3.5cm]
    at (5.5, -3.5) {\texttt{KCONFIG\_SYSCALL\_BACKEND}\\[2pt]0 = int0x80, 1 = sysenter};

\end{tikzpicture}

  \caption{\texttt{syscall\_ops\_t} 可插拔后端架构——后端负责入口，dispatch 负责分发}
  \label{fig:syscall-dispatch}
\end{figure}

\begin{designbox}
为什么 \texttt{syscall\_ops\_t} 只有一个 \texttt{init} 函数指针？

与 \texttt{vma\_ops\_t}（7 个函数指针）不同，系统调用后端只需做一件事：
\textbf{设置入口点}。一旦入口安装完毕，后续所有系统调用都走同一条路径：
用户态 → 汇编 stub → \texttt{syscall\_dispatch()} → handler。

后端不需要 \texttt{dispatch} 或 \texttt{handle} 函数指针，
因为汇编 stub 直接调用 dispatch 层的 C 函数。
将来如果需要后端特定的 teardown（例如 MSR 清理），
可以添加 \texttt{fini} 函数指针——但现在遵循 YAGNI 原则，不提前设计。
\end{designbox}

% ----------------------------------------------------------------------------
\subsection{后端注册 API}
% ----------------------------------------------------------------------------

\codefile{src/include/kernel/syscall.h}
\begin{lstlisting}[style=cstyle]
/* 注册后端（在 syscall_init 之前调用） */
void syscall_register_backend(const syscall_ops_t* ops);

/* 返回当前后端名称，NULL=未注册 */
const char* syscall_backend_name(void);

/* 系统调用子系统是否已初始化 */
int syscall_is_ready(void);
\end{lstlisting}

% ----------------------------------------------------------------------------
\subsection{handler 函数签名}
% ----------------------------------------------------------------------------

每个具体的系统调用（如 \texttt{sys\_write}）实现为此签名的函数，
注册到 syscall table 中：

\begin{lstlisting}[style=cstyle]
typedef int32_t (*syscall_handler_t)(const syscall_regs_t* regs);
\end{lstlisting}

\begin{itemize}
  \item 参数：\texttt{syscall\_regs\_t} 包含了调用号和最多 5 个参数
  \item 返回值：由后端写回用户态 EAX。负数表示 $-$errno（如 $-38 = -$ENOSYS）
\end{itemize}

% ----------------------------------------------------------------------------
\subsection{注册单个 handler}
% ----------------------------------------------------------------------------

\begin{lstlisting}[style=cstyle]
/*
 * syscall_register -- 注册/替换一个 syscall handler
 *
 * @param nr       系统调用号（0 < nr < SYSCALL_MAX）
 * @param handler  处理函数
 * @return         0=成功, -1=nr 越界
 */
int syscall_register(uint32_t nr, syscall_handler_t handler);
\end{lstlisting}

这个 API 不仅供内核初始化使用，将来用户态驱动（微内核架构）也可以
通过系统调用注册新的 handler——实现内核服务的动态扩展。
